% TOP_PROBLEM: {"id": 2590, "title": "Kourovka Notebook Problem 21.81", "category_id": 17, "category": {"id": 17, "name": "group_theory", "display_name": "Group Theory", "description": "Problems about groups, group actions, representations, and related algebraic structures.", "slug": "group-theory", "order_index": 17, "created_at": "2026-07-31T15:26:25.671Z"}, "status": "open", "problem_number": "KOU-21.81", "set_id": 10}
% FIRST_POSTED: 2026-10-01
% ENTRY_KIND: statement_only
% UnsolvedMath ID 2590; source catalogue code KOU-21.81
% !TeX program = lualatex
% Definition and sources only; this file is not an LLM solution attempt.
\documentclass[11pt,a4paper]{article}
\usepackage[margin=25mm]{geometry}
\usepackage{fontspec}
\setmainfont{lmroman10-regular.otf}[BoldFont=lmroman10-bold.otf,
 ItalicFont=lmroman10-italic.otf,BoldItalicFont=lmroman10-bolditalic.otf]
\usepackage{amsmath,amssymb,mathtools}
\usepackage{unicode-math}
\setmathfont{latinmodern-math.otf}
\AtBeginDocument{\let\setminus\smallsetminus}
\usepackage{xurl,hyperref}
\hypersetup{colorlinks=true,urlcolor=blue}
\setcounter{secnumdepth}{0}
\setlength{\parindent}{0pt}
\setlength{\parskip}{5pt}
\setlength{\emergencystretch}{3em}
\begin{document}
\section*{Kourovka Notebook Problem 21.81}\label{2590}
{\small UnsolvedMath ID 2590 · KOU-21.81 · Group Theory · Kourovka Notebook - New Problems, Issue 21}
\subsection{Definitions and mathematical statement}
For an integer $n\ge3$, let $F_n=\langle x_1,\ldots,x_n\rangle$ be the free group on $n$ generators: arbitrary choices of $n$ elements of any group extend uniquely to a homomorphism from $F_n$. Let $\Gamma$ be a finite simple group, meaning a nontrivial finite group with no normal subgroups other than $1$ and itself. Define
\[
 \mathcal N_n(\Gamma)=\{N\trianglelefteq F_n:F_n/N\cong\Gamma\}.
\]
These are exactly the kernels of surjective homomorphisms $F_n\to\Gamma$. The automorphism group $\operatorname{Aut}(F_n)$ acts on this set by sending $N$ to $\alpha(N)$.

The conjecture is that this action is transitive for every such $n$ and $\Gamma$. Explicitly, for any $N_1,N_2\in\mathcal N_n(\Gamma)$, there should exist $\alpha\in\operatorname{Aut}(F_n)$ with $\alpha(N_1)=N_2$.

Equivalently, given any two epimorphisms $\varphi,\psi:F_n\to\Gamma$, there should be $\alpha\in\operatorname{Aut}(F_n)$ and $\beta\in\operatorname{Aut}(\Gamma)$ such that $\psi=\beta\circ\varphi\circ\alpha$. Epimorphisms with the same kernel differ by an automorphism of their common quotient, explaining the appearance of $\beta$. Accordingly the proposed transitivity is on quotient kernels, with no chosen identification of the quotient with $\Gamma$. The stated range begins at rank three and imposes no upper bound on the order of $\Gamma$.
\subsection{Short English statement}
For every finite simple group, are all its quotient kernels in a free group of rank at least three equivalent under free-group automorphisms?
\subsection{Sources}
\begin{itemize}
\item[S1] ULAM.AI and the UnsolvedMath Contributors, supplied problems.json, record 2590 (KOU-21.81), Kourovka Notebook - New Problems, Issue 21. Catalogue identity and source transcription; CC BY 4.0.
\url{https://huggingface.co/datasets/ulamai/UnsolvedMath}.
\item[S2] The Kourovka Notebook, 21st edition, version 47 (30 September 2026), new problems, Problem 21.81. Expanded formulation of the supplied catalogue statement.
\url{https://arxiv.org/pdf/1401.0300v47}.
\end{itemize}
\end{document}
